% Web専用原稿．
% 構成案：両側イデアル・生成・商環と同型定理．
% 文献確認：Lang, Algebra, Ch. II（イデアルと剰余環の本文を確認）．
\section{イデアル}
% genron:begin K-000327
% genron:use K-000001
% genron:use K-000002
% genron:use K-000003
% genron:use K-000004
% genron:use K-000005
% genron:use K-000007
% genron:use K-000017
% genron:use K-000249
% genron:use K-000228
\begin{defn}[イデアル]
環 $A$ の加法群の部分群 $\mathfrak a$ が任意の $a\in A$，$x\in\mathfrak a$ に対して $ax,xa\in\mathfrak a$ を満たすとき，$\mathfrak a$ を（両側）イデアルという．可換環では左右いずれかの条件で足りる．
\end{defn}
% genron:end K-000327
% genron:begin K-000328
% genron:use K-000001
% genron:use K-000002
% genron:use K-000003
% genron:use K-000004
% genron:use K-000005
% genron:use K-000007
% genron:use K-000017
% genron:use K-000327
% genron:use K-000230
% genron:use K-000250
% genron:use K-000231
% genron:use K-000165
\begin{prop}[生成イデアル]
$A$ のイデアルの任意の共通部分はイデアルである（空の共通部分は $A$ とする）．部分集合 $S\subseteq A$ を含む最小のイデアル $(S)$ は，$asb$（$a,b\in A$，$s\in S$）の有限和全体である．
\end{prop}
\begin{proof}
共通部分は加法部分群であり，左右からの積について閉じる．有限和全体は空和 $0$ を含み，和と逆元及び左右からの積で閉じる（逆元は $a$ を $-a$ に替える）．$1s1=s$ であるから $S$ を含む．$S$ を含む任意のイデアルは各 $asb$ を含み，有限和も含むので最小性を得る．
\end{proof}
% genron:end K-000328
% genron:begin K-000329
% genron:use K-000001
% genron:use K-000002
% genron:use K-000003
% genron:use K-000004
% genron:use K-000005
% genron:use K-000007
% genron:use K-000017
% genron:use K-000327
% genron:use K-000251
% genron:use K-000285
% genron:use K-000286
% genron:use K-000287
\begin{prop}[環準同型の核]
環準同型 $f:A\to B$ の核 $\ker f=\{a\in A\mid f(a)=0\}$ は $A$ のイデアルである．$B$ のイデアルの逆像はイデアルであり，$f$ が全射なら $A$ のイデアルの像もイデアルである．
\end{prop}
\begin{proof}
加法部分群に関する主張は \GenronRef{K-000286} による．逆像について $f(ax)=f(a)f(x)$，$f(xa)=f(x)f(a)$ であるから吸収性が従う．核は零イデアルの逆像である．像については $b\in B$ を $f(a)$ と書くことができ，$bf(x)=f(ax)$，$f(x)b=f(xa)$ を用いる．
\end{proof}
% genron:end K-000329
\section{商環}
% genron:begin K-000330
% genron:use K-000001
% genron:use K-000002
% genron:use K-000003
% genron:use K-000004
% genron:use K-000005
% genron:use K-000007
% genron:use K-000017
% genron:use K-000327
% genron:use K-000267
% genron:use K-000249
% genron:use K-000250
% genron:use K-000268
\begin{prop}[商環]
イデアル $\mathfrak a$ に対し，加法商群 $A/\mathfrak a$ 上で
\[
(a+\mathfrak a)(b+\mathfrak a)=ab+\mathfrak a
\]
と定めると環となり，標準全射 $q:A\to A/\mathfrak a$ は環準同型である．
\end{prop}
\begin{proof}
加法群は可換であるから $\mathfrak a$ は正規部分群である．$a'=a+x$，$b'=b+y$（$x,y\in\mathfrak a$）なら
$a'b'-ab=ay+xb+xy\in\mathfrak a$ であり，積は代表元に依存しない．乗法の結合律・単位律・分配律は代表元についての環の法則から従う．標準全射は各演算と単位元を保存する．
\end{proof}
% genron:end K-000330
% genron:begin K-000331
% genron:use K-000001
% genron:use K-000002
% genron:use K-000003
% genron:use K-000004
% genron:use K-000005
% genron:use K-000007
% genron:use K-000017
% genron:use K-000330
% genron:use K-000288
% genron:use K-000251
\begin{prop}[商環の普遍性]
イデアル $\mathfrak a$ と環準同型 $f:A\to B$ に対し，$f=\bar f\circ q$ を満たす環準同型 $\bar f:A/\mathfrak a\to B$ が存在することと $\mathfrak a\subseteq\ker f$ は同値である．存在するとき $\bar f$ は一意である．
\end{prop}
\begin{proof}
必要性は $q(x)=0$（$x\in\mathfrak a$）による．十分性については \GenronRef{K-000288} を加法群に適用し，$\bar f(a+\mathfrak a)=f(a)$ と定まる．代表元で積と単位元を比較すれば環準同型である．一意性は $q$ の全射性による．
\end{proof}
% genron:end K-000331
% genron:begin K-000332
% genron:use K-000001
% genron:use K-000002
% genron:use K-000003
% genron:use K-000004
% genron:use K-000005
% genron:use K-000007
% genron:use K-000017
% genron:use K-000331
% genron:use K-000289
% genron:use K-000252
\begin{prop}[環準同型の第一同型定理]
環準同型 $f:A\to B$ に対して
\[
 A/\ker f\longrightarrow f(A),\qquad a+\ker f\longmapsto f(a)
\]
は環同型である．
\end{prop}
\begin{proof}
前命題によって環準同型が定まり，\GenronRef{K-000289} を加法群に適用すると全単射である．\GenronRef{K-000252} により逆写像も環準同型である．
\end{proof}
% genron:end K-000332
% genron:begin K-000333
% genron:use K-000329
% genron:use K-000001
% genron:use K-000002
% genron:use K-000003
% genron:use K-000004
% genron:use K-000005
% genron:use K-000007
% genron:use K-000017
% genron:use K-000330
% genron:use K-000292
% genron:use K-000327
% genron:use K-000251
% genron:use K-000286
\begin{prop}[イデアル対応]
標準全射 $q:A\to A/\mathfrak a$ の像と逆像は，$\mathfrak a$ を含む $A$ のイデアル全体と $A/\mathfrak a$ のイデアル全体との間の互いに逆な順序保存写像を定める．
\end{prop}
\begin{proof}
像及び逆像がイデアルとなることは\GenronRef{K-000329}と同じ吸収性の論証による．\GenronRef{K-000292} により加法部分群について像と逆像は互いに逆であり，包含を保存する．
\end{proof}
% genron:end K-000333
% genron:begin K-000334
% genron:use K-000001
% genron:use K-000002
% genron:use K-000003
% genron:use K-000004
% genron:use K-000005
% genron:use K-000007
% genron:use K-000017
% genron:use K-000327
\begin{prop}[単位元を含むイデアル]
イデアル $\mathfrak a$ が $1_A$ を含めば $\mathfrak a=A$ である．従って真のイデアルは単位元を含まない．
\end{prop}
\begin{proof}
任意の $a\in A$ は $a1_A\in\mathfrak a$ であるから $\mathfrak a=A$ である．
\end{proof}
% genron:end K-000334
